\documentclass[12pt,a4paper]{article}
\usepackage[T2A]{fontenc}
\usepackage[utf8]{inputenc}
\usepackage[russian]{babel}
\usepackage{amsmath,amssymb}
\usepackage{geometry}
\usepackage{graphicx}
\usepackage{float}
\usepackage{booktabs}
\usepackage{array}
\usepackage{longtable}
\usepackage{caption}
\usepackage{hyperref}
\geometry{margin=2.2cm}

\title{Сведение задачи оптимального управления к задаче минимизации функции многих переменных}
\author{Буканов Иван, группа 402}
\date{Апрель 2026}

\begin{document}

\newpage
\maketitle
\newpage

\section{Постановка задачи}

Рассматривается задача оптимального управления с терминальным функционалом
\[
    J = y(t_1) \to \min
\]
при ограничениях
\[
    \dot x = u, \qquad \dot y = x^2,
\]
\[
    -1 \le u(t) \le 1,
\]
\[
    x(t_0)=x_0, \qquad x(t_1)=x_1, \qquad y(t_0)=0.
\]

Переменная \(y\) используется для перехода от интегрального функционала к терминальному. Так как
\[
    \dot y = x^2,
\]
то значение \(y(t_1)\) равно накопленному интегралу
\[
    y(t_1)=\int_{t_0}^{t_1} x^2(t)\,dt.
\]
Следовательно, минимизация \(y(t_1)\) означает, что траектория должна по возможности быстрее приближаться к нулю, а затем перейти в заданное конечное состояние \(x_1\).

\section{Дискретизация}

Разобьём отрезок \([t_0,t_1]\) на \(N\) равных частей:
\[
    t_i = t_0 + i\,dt, \qquad i=0,\ldots,N,
\]
где
\[
    dt = \frac{t_1-t_0}{N}.
\]

В качестве переменных конечномерной задачи берём значения
\[
    x_0,\ldots,x_N, \qquad
    y_0,\ldots,y_N, \qquad
    u_0,\ldots,u_{N-1}.
\]
Все эти величины объединяются в один вектор
\[
    z = (x_0,\ldots,x_N,\; y_0,\ldots,y_N,\; u_0,\ldots,u_{N-1}).
\]

Для уравнений движения используется явная схема Эйлера:
\[
    x_{i+1} = x_i + dt\,u_i,
    \qquad i=0,\ldots,N-1,
\]
\[
    y_{i+1} = y_i + dt\,x_i^2,
    \qquad i=0,\ldots,N-1.
\]

После дискретизации исходная задача сводится к задаче минимизации функции многих переменных при ограничениях на динамику, граничные условия и управление.

\section{Метод штрафных функций}

Ограничения учитываются с помощью штрафа. Для фиксированного параметра \(\rho>0\) минимизируется функция
\[
    \Phi_\rho(z) = y_N + \rho P(z),
\]
где \(P(z)\) --- сумма квадратов нарушений всех ограничений.

В работе использован штраф
\[
\begin{aligned}
P(z) ={}&
    \sum_{i=0}^{N-1}
    \left(x_{i+1}-x_i-dt\,u_i\right)^2
    +
    \sum_{i=0}^{N-1}
    \left(y_{i+1}-y_i-dt\,x_i^2\right)^2
    \\
    &+
    (x_0-\bar x_0)^2
    +
    (x_N-\bar x_1)^2
    +
    y_0^2
    \\
    &+
    \sum_{i=0}^{N-1}
    \left(\max(0,u_i-1)^2+\max(0,-1-u_i)^2\right).
\end{aligned}
\]

Здесь \(\bar x_0\) и \(\bar x_1\) --- заданные значения левого и правого граничных условий. Последняя сумма отвечает за выполнение ограничения
\[
    -1 \le u_i \le 1.
\]

Для усиления выполнения ограничений задача решается несколько раз подряд при возрастающих значениях \(\rho\). Найденное решение для предыдущего значения \(\rho\) используется как начальное приближение для следующего.

\section{Численная реализация}

В ноутбуке использованы следующие параметры:
\[
    t_0 = 0, \qquad t_1 = 4,
\]
\[
    \bar x_0 = 2, \qquad \bar x_1 = -1,
\]
\[
    N = 60.
\]

Начальное приближение выбирается следующим образом:
\begin{itemize}
    \item значения \(x_i\) задаются линейной интерполяцией между \(\bar x_0\) и \(\bar x_1\);
    \item значения \(u_i\) вычисляются по разностям соседних значений \(x_i\);
    \item значения \(y_i\) вычисляются накоплением величин \(dt\,x_i^2\).
\end{itemize}

Минимизация функции \(\Phi_\rho\) проводится методом L-BFGS-B. В ноутбуке используется последовательность
\[
    \rho = 1,\;10,\;100,\;1000,\;10000.
\]

После каждого шага сохраняются следующие контрольные величины:
\begin{itemize}
    \item значение функционала \(J=y_N\);
    \item максимальная невязка уравнения для \(x\);
    \item максимальная невязка уравнения для \(y\);
    \item максимальное нарушение ограничения на управление.
\end{itemize}

\section{Результаты}

После выполнения минимизации получаются массивы \(x_i\), \(y_i\) и \(u_i\). По ним строятся точечные графики, которые позволяют проверить характер найденного управления и траектории.

\begin{figure}[H]
    \centering
    \includegraphics[width=0.85\textwidth]{materials/x_plot.png}
    \caption{Точечный график \(x(t_i)\).}
\end{figure}

\begin{figure}[H]
    \centering
    \includegraphics[width=0.85\textwidth]{materials/y_plot.png}
    \caption{Точечный график \(y(t_i)\).}
\end{figure}

\begin{figure}[H]
    \centering
    \includegraphics[width=0.85\textwidth]{materials/u_plot.png}
    \caption{Точечный график \(u(t_i)\).}
\end{figure}

\section{Вывод}

Исходная задача оптимального управления была сведена к задаче минимизации функции многих переменных. Для этого была введена равномерная сетка, уравнения движения были заменены явной схемой Эйлера, а ограничения были включены в целевую функцию с помощью штрафных слагаемых.

Ожидаемый вид решения соответствует смыслу задачи: траектория \(x(t)\) стремится как можно быстрее приблизиться к нулю, чтобы уменьшить интеграл от \(x^2\), после чего переходит к заданному конечному значению \(x_1\). Управление при этом должно оставаться в допустимом диапазоне \([-1,1]\).

\end{document}
